\todo{Draft; verify every citation before submission.}

\paragraph{Activation steering.} Adding a fixed vector to the residual stream at inference time changes
behaviour without changing weights: ActAdd \citep{turner2023actadd}, contrastive activation addition
\citep{panickssery2023caa}, inference-time intervention \citep{li2023iti}. Refusal in chat models is mediated
by a single direction that can be removed or added \citep{arditi2024refusal}. Our intervention is a special
case: the steering vector is one column of $\Wdown$, so the direction is chosen from the model's own write
vectors rather than fitted.

\paragraph{Neurons as features.} Feed-forward layers act as key-value memories whose values write
interpretable directions \citep{geva2021kv}; individual neurons have been linked to knowledge
\citep{dai2022knowledge} and to context features \citep{gurnee2023haystack}. Sparse autoencoders recover
many more monosemantic features than there are neurons \citep{templeton2024scaling}, which makes it
surprising that single raw neurons suffice here.

\paragraph{Persona and character.} Persona vectors \citep{chen2025persona} are residual directions fitted to
trait contrasts and used to monitor or steer traits; emergent misalignment \citep{betley2025emergent} shows
that narrow fine-tuning can flip a model's whole persona. Our result sits between them: the direction is a
single neuron's write vector, found by cosine against a persona-vector-like contrast.

\paragraph{The model.} Qwen3.8-27B has 64 decoder layers, a 5{,}120-wide residual stream and a
17{,}408-wide MLP; three of every four layers use Gated DeltaNet linear attention
\citep{qwen2026qwen38,yang2024gated}. Steering sites 33 and 34 are linear-attention layers.
